\documentclass[11pt]{article}
\usepackage[margin=1in]{geometry}
\usepackage{titlesec}
\usepackage{hyperref}
\usepackage[numbers]{natbib}
\usepackage{amsmath}
\usepackage{amssymb}
\usepackage{graphicx}
\usepackage{booktabs}
\usepackage{setspace}
\usepackage{xcolor}
\usepackage{enumitem}
\usepackage{parskip}

\hypersetup{
    colorlinks=true,
    linkcolor=blue!60!black,
    citecolor=blue!60!black,
    urlcolor=blue!60!black
}

\titleformat{\section}{\large\bfseries}{\thesection.}{0.5em}{}
\titleformat{\subsection}{\normalsize\bfseries}{\thesubsection.}{0.5em}{}
\titlespacing*{\section}{0pt}{1.2ex plus 0.5ex minus 0.2ex}{0.8ex plus 0.2ex}
\titlespacing*{\subsection}{0pt}{1.0ex plus 0.5ex minus 0.2ex}{0.5ex plus 0.2ex}

\title{
    \vspace{-1.5cm}
    \textbf{Can adding Assymetric Relation Dynamics make Knowledge Transfer more Efficient in Language Agents?} \\[0.4em]
    \large Project Proposal - EN.601.773 Machine Social Intelligence, Spring 2026
}
\author{Rishi More (rmore2)}
\date{\today}

\begin{document}
\maketitle
\vspace{-0.5cm}
% \hrule
\vspace{0.4cm}


% -------------------------------------------------------
\section{Motivation and Research Question}

Human intelligence does not develop in isolation. From birth, teaching and learning are performed through a \emph{relationship}: the mother (caregiver) knows more, senses more, and performs actions to selectively reveal information that the child (learner) is ready to receive. What the infant absorbs is not a random snapshot of the world but a curated stream of experience, decided by the caregiver's theory of the child's mind \citep{gweon2021inferential,csibra2009natural}. 

This teamwork between an expert agent and a developing one should be the most efficient knowledge-transfer mechanism in nature, but it is absent from how AI agents are trained nowadays.

Through this project, I want to investigate the question:\\
\textbf{Does learning in an asymmetric relational dynamic (such as a ``mother--child'' structure) enable an agent to acquire richer and more generalizable knowledge than learning alone or from a symmetric peer?}

I plan to implement this by equipping agents with a memory (like human memory) architecture, comprising instincts, short-term memory, long-term memory, and habits, and examining how relational context shapes what is stored, compressed, and transferred.

Crucially, I plan to design the framework to run within a text-only, LoRA fine-tuning pipeline, making the project fully self-contained without external simulation environments. Therefore, all memory modules will be implemented as lightweight Python orchestration around Tinker's primitives.

% -------------------------------------------------------
\section{Background and Related Work}

\paragraph{Memory architecture in cognitive agents.}
 Neuroscience distinguishes at least four interacting systems relevant here \citep{tulving1985memory}: 
 \begin{enumerate}
     \item Instinctive memory encoding primitive behavioral priors.
     \item Working (short-term) memory holding a limited and rapidly decaying context
     \item Long-term memory storing temporally indexed personal experiences
     \item Habits comprising stimulus--response associations created by repetition that bypass normal thinking and planning. 
 \end{enumerate}
 
 Current computational agents tend to reflect a somewhat similar approach: \citet{park2023generative} uses a hierarchical memory stream with reflection, retrieval, and planning, but doesn't examine how social relationships shape memory formation/modification.

\paragraph{Social learning and pedagogy.}
Natural pedagogy theory \citep{csibra2009natural} says that humans are uniquely adapted to both teach and be taught: caregivers \emph{ostensively} signal when information is generic vs. worth storing, while learners are inclined to learn more from information delivered in pedagogical frames. \citet{gweon2021inferential} formalizes this as a cooperative inference process in which the teacher models the learner's knowledge state to select demonstrations, and the learner uses the teacher's selections as evidence to update beliefs.

This modeling (in theory) produces faster and more accurate learning than passive observation. Machine social learning work \citep{ho2019cooperative} has explored cooperative IRL between humans and agents, but has not examined whether a persistent relationship between agents is itself informative beyond the individual teaching actions.

\paragraph{Relational context as an inductive bias.}
There is growing evidence that social context acts as a bias on what gets encoded. Emotionally significant or relationally relevant events should technically occupy more space in the memory than other events \citep{dolcos2004interaction}. 

Prior work on LLMs with augmented personalities \citep{shao2023character} shows that adding stable roles to LLMs shifts response distributions, but this has not been studied specifically for memory formation or learning speed.

\paragraph{LoRA fine-tuning as habit formation.}
Recent work on parameter-efficient fine-tuning \citep{hu2022lora} demonstrates that low-rank weight adaptations can encode behavioral biases that persist across contexts. 

Therefore, I plan to use this LoRA fine-tuning to approximate the habit-formation mechanism. Repeated exposure to similar action-response patterns in the relational setting should ideally produce stronger adaptations than the same exposure in a non-relational setting.

% -------------------------------------------------------
\section{Proposed Framework}

\subsection{Text-Only Household Task Environment}

Rather than relying on external simulation platforms, I want to construct a text-based household task benchmark consisting of \~200 multi-step procedural scenarios (e.g., preparing a meal, doing laundry, assembling furniture). Each scenario is represented as a structured natural-language problem: a goal description, a set of available objects, and a sequence of actions required to complete the task. The caregiver has access to the complete action sequence and affordance map; the child does not. Tasks are partitioned into a training curriculum of 160 scenarios (graded from 1-step to 8-step tasks) and a held-out transfer set of 40 novel scenarios used exclusively for evaluation. This design keeps the entire pipeline within text, making it fully compatible with Tinker's interface.

\subsection{Agent Memory Architecture}

Each agent has with four memory components implemented as Python objects wrapped around Tinker API calls:

\begin{enumerate}[leftmargin=*, label=\textbf{M\arabic*.}]
    \item \textbf{Instinct Buffer} ($\mathcal{I}$): A fixed system prompt that encodes role-specific behavioral priors and is \emph{never updated} during training. 
    
    For the caregiver: 
    \emph{``You are a knowledgeable parent. Your goal is to help your child understand how to complete tasks by explaining your reasoning and filling in their knowledge gaps.''} 
    
    For the child: \emph{``You are a curious child learning about the world. Ask questions and try to understand why things work the way they do.''} 
    
    The instinct buffer is passed as the system prompt in every \texttt{sample()} call for its agent but is excluded from the LoRA training objective, preserving it as a true fixed prior.

    \item \textbf{Short-Term Working Memory} ($\mathcal{W}_t$): A sliding window of the last $k{=}8$ dialogue turns, passed directly as the conversation history in each \texttt{sample()} call. Older entries are removed from the context when the window is full.

    \item \textbf{Long-Term Memory} ($\mathcal{L}$): A plain-text episodic store maintained as a Python list of summarized past episodes. After each episode, a lightweight summarization step (one additional call with a compression prompt) compresses the episode into a 2-3 sentence structured entry recording: the task, what the child did not know, what the caregiver explained, and whether the task succeeded. Retrieval at the start of each new episode selects the top-3 most relevant past entries by simple keyword overlap (BM25) with the new task description and adds them to $\mathcal{W}_t$. 
    
    (*I deliberately avoid vector databases and embedding models to keep the implementation dependency-free.)

    \item \textbf{Habit Schema Store} ($\mathcal{H}$): Implemented directly as the LoRA adapter weights. When an episode is marked as a successful habit creation candidate (defined as the child producing the correct action sequence with no caregiver correction for the third consecutive time on a given task), the child agent's \texttt{forward\_backward()} is called on the successful trajectory and \texttt{optim\_step()} is applied. 
    
    (The LoRA weights should accumulate the agent's habituated responses over training, acting like a persistent behavioral cache that does not require explicit retrieval.)
\end{enumerate}

\paragraph{Salience signal and consolidation gate.}
The central hypothesis is that relational context controls the remember signal (formally called the ``salience signal``), thus controlling what gets written to $\mathcal{L}$ and triggers habit formation in $\mathcal{H}$. 

Salience score can be defined as (for an episode):
\[
s = \alpha \cdot \text{novelty}(e) + \beta \cdot \text{emotion}(e) + \gamma \cdot \text{teaching}(e)
\]
where $\text{novelty}(e)$ is 1 if the task class has not appeared in $\mathcal{L}$; $\text{emotion}(e) \in [-1,1]$ is a sentiment score over the child's utterances (computed via a frozen small classifier); and $\text{teaching}(e) \in \{0,1\}$ flags whether the caregiver produced at least one explicitly corrective or explanatory utterance. 

An episode is added into $\mathcal{L}$ if $s > \tau$ (threshold for long-term memory). The scalars $\alpha, \beta, \gamma$ are simple hyperparameters. In the non-relational baseline, $\gamma \equiv 0$ (there is no teacher).

\subsection{Interaction Protocol: The Mother--Son Dyad}
Plan to simulate asymmetric role assignment with two Tinker-hosted model instances:

\begin{itemize}
    \item \textbf{Caregiver agent ($A_C$, ``Mother''):} Initialized from a Qwen3-8B-Instruct checkpoint with a caregiver system prompt ($\mathcal{I}_C$), has access to the full task solution. At each turn, $A_C$ calls \texttt{sample()} to generate either a targeted hint, a corrective explanation, or a question probing the child's understanding. 
    
    $A_C$'s weights are \emph{frozen} throughout training; it acts as a fixed expert teacher (although it seems interesting to observe how the teacher learns to teach through teaching episodes, looks too ambitious). 
    
    This eliminates the need to train two agents simultaneously and keeps compute costs within the \$249 budget.

    \item \textbf{Learner agent ($A_L$, ``Child''):} Initialized from a Qwen3-8B base checkpoint with a child system prompt ($\mathcal{I}_L$), has no access to the task solution. At each turn, $A_L$ calls \texttt{sample()} to generate an action attempt or a question. 
    
    $A_L$'s LoRA adapter is updated via \texttt{forward\_backward()} and \texttt{optim\_step()} on successful habit-forming episodes.
\end{itemize}

A single episode proceeds as follows: 
\begin{enumerate}
    \item Task description is presented to both agents
    \item Child ($A_L$) attempts a step
    \item Mother ($A_C$) observes and responds with a hint or confirmation
    \item Repeat steps 2 and 3 until success or a 15-turn limit
    \item Reward is computed
    \item Long-term memory and LoRA are updated using the salience score 
    \end{enumerate}

\subsection{Experimental Conditions}

To isolate the effect of relational structure, I will compare four conditions using the same base model (Qwen3-8B):

\begin{center}
\begin{tabular}{p{2.8cm}p{4.5cm}p{5.5cm}}
\toprule
\textbf{Condition} & \textbf{Agent Structure} & \textbf{Memory \& Training} \\
\midrule
Solo & Single agent, no teacher & $\mathcal{W} + \mathcal{L}^{\text{self}} + \mathcal{H}^{\text{LoRA}}$; memory augmentation only by task success \\
Symmetric Peer & Two identical agents, no role asymmetry & $\mathcal{W} + \mathcal{L} + \mathcal{H}^{\text{LoRA}}$ (both); $\gamma \equiv 0$ \\
Role-Labeled & Asymmetric prompts, no teaching salience & $\mathcal{I} + \mathcal{W} + \mathcal{L} + \mathcal{H}^{\text{LoRA}}$; $\gamma \equiv 0$ \\
\textbf{Relational (Mine)} & Asymmetric roles + relational salience score & $\mathcal{I} + \mathcal{W} + \mathcal{L}^{rel} + \mathcal{H}^{\text{LoRA}}$; full $s$ with $\gamma > 0$ \\
\bottomrule
\end{tabular}
\end{center}

(The Role-Labeled condition controls for whether prompt asymmetry alone, without the teaching salience signal, accounts for any observed advantage.)

% -------------------------------------------------------
\section{Research Hypotheses}

\begin{enumerate}[label=\textbf{H\arabic*.}]
    \item \textbf{Better memory:} The relational condition will create a child agent whose long-term memory contains more accurate and generalizable entries (measured by transfer task accuracy) after the same number of training episodes.

    \item \textbf{Habit acceleration:} LoRA habits will be formed in fewer episodes in the Relational condition than in Solo or Peer, because the mother's help creates repeated, consistently correct trajectories.

    \item \textbf{ToM proxy in child:} After training, the child agent's ability to predict which hint the caregiver will give next (measured on a held-out dialogue prediction task) should be higher in the Relational condition, due to possible internalization of the caregiver's pedagogical model.
\end{enumerate}

\paragraph{Budget estimate:}
At Qwen3-8B pricing (\$0.40/M sample tokens, \$0.40/M train tokens), running 160 training episodes $\times$ 15 turns $\times$ 2 agents $\times$ $\sim$500 tokens/turn yields approximately 2.4M sample tokens per condition, or $\sim$\$1.00 per condition. Four conditions $\times$ 3 seeds $\approx$ \$12 for training sampling. LoRA updates are triggered sparsely (estimated 30\% of episodes), adding $\sim$\$3. Evaluation on 40 held-out tasks adds $\sim$\$2. \textbf{Total estimated cost: under \$20}.

% -------------------------------------------------------
\bibliographystyle{unsrtnat}
\bibliography{bibliography}

\end{document}